% Assignment source: ne630/homework/html/HW13.html.
% Legacy foundation: ne630_problems/lesson_13/statement.tex and hw13_sol.ipynb.
% Problem 1 restores the part (b) calculation that was commented in the legacy
% statement source.

\noindent{\bf Problem 1}. One can compute effective cross sections for a
coarse group structure based on the multigroup fluxes and cross sections
defined on a finer group structure. Let a fine energy group $g$ be defined over
the energies $E_g<E\le E_{g-1}$. Suppose a coarse group $G$ is defined over
$E_{g+1}<E\le E_{g-1}$ (i.e., over groups $g$ and $g+1$). The coarse-group
flux can then be written as
\[
\begin{aligned}
 \phi_G^{\mathrm{coarse}}
 &=\int_{E_{g+1}}^{E_{g-1}}\phi(E)\,dE\\
 &=\int_{E_g}^{E_{g-1}}\phi(E)\,dE
  +\int_{E_{g+1}}^{E_g}\phi(E)\,dE\\
 &=\phi_g^{\mathrm{fine}}+\phi_{g+1}^{\mathrm{fine}}.
\end{aligned}
\]
\begin{enumerate}[label=(\alph*)]
\item Determine an expression for $\sigma_{a,G}$ in terms of the fine-group
      fluxes and cross-section values $\sigma_{a,g}$ and
      $\sigma_{a,g+1}$.
\item Assuming the fast and epithermal fluxes are identical, determine the
      effective $\sigma_a$ for \isoA{U}{238} over
      $1<E<10^7~\mathrm{eV}$ using FNRP Table 3.2.
\end{enumerate}

\begin{adaptedsolution}
The reaction rate per atom integrated over the coarse group must be preserved:
\[
 \phi_G^{\mathrm{coarse}}\sigma_{a,G}
 =\phi_g^{\mathrm{fine}}\sigma_{a,g}
 +\phi_{g+1}^{\mathrm{fine}}\sigma_{a,g+1}.
\]
Therefore,
\[
 \boxed{\sigma_{a,G}
 =\frac{\phi_g^{\mathrm{fine}}\sigma_{a,g}
       +\phi_{g+1}^{\mathrm{fine}}\sigma_{a,g+1}}
      {\phi_g^{\mathrm{fine}}+\phi_{g+1}^{\mathrm{fine}}}}.
\]
For part (b), Table~3.2 gives $I_a\approx278~\mathrm b$ for the epithermal
group, so
\[
 \sigma_a^{\mathrm{epi}}
 =0.0869 I_a
 \approx24.16~\mathrm b.
\]
The fast-group value is $\sigma_a^{\mathrm{fast}}=0.361~\mathrm b$.  With
identical fast and epithermal fluxes,
\[
 \boxed{\bar{\sigma}_a
 =\frac{24.16+0.361}{2}
 \approx12.26~\mathrm b}.
\]
\end{adaptedsolution}

\newpage

\noindent{\bf Problem 2}. Consider an infinite solution of water and uranium
with a uniform source of $1~\mathrm{MeV}$ neutrons that emits
$100~\mathrm{n\,cm^{-3}\,s^{-1}}$. Using the data in Table 1, compute the
volumetric absorption and fission rates in the solution.

\begin{center}
\textbf{Table 1.} Two-group data for Problems 2 and 3. Macroscopic cross
sections are given in $\mathrm{cm^{-1}}$.
\[
\begin{array}{lrr}
\toprule
 & g=1 & g=2\\
\midrule
\Sigma_\gamma & 2.04\times10^{-3} & 1.91\times10^{-2}\\
\Sigma_f & 2.96\times10^{-4} & 1.38\times10^{-2}\\
\bar{\nu}\Sigma_f & 7.38\times10^{-4} & 3.36\times10^{-2}\\
\Sigma_{s,1\gets g} & 7.21\times10^{-1} & 0\\
\Sigma_{s,2\gets g} & 4.27\times10^{-2} & 2.72\times10^0\\
\bottomrule
\end{array}
\]
\end{center}

\begin{adaptedsolution}
The two-group equations for this fixed-source infinite system are
\begin{align*}
 \Sigma_{r,1}\phi_1
 &=\bar{\nu}\Sigma_{f,1}\phi_1
  +\bar{\nu}\Sigma_{f,2}\phi_2+S_1,\\
 \Sigma_{a,2}\phi_2&=\Sigma_{s,2\gets1}\phi_1.
\end{align*}
Thus,
\[
 \phi_2=\frac{\Sigma_{s,2\gets1}\phi_1}{\Sigma_{a,2}},
 \qquad
 \phi_1=
 \frac{S_1}
 {\Sigma_{r,1}-\bar{\nu}\Sigma_{f,1}
 -\bar{\nu}\Sigma_{f,2}\Sigma_{s,2\gets1}/\Sigma_{a,2}}.
\]
The fast removal cross section is
\[
 \Sigma_{r,1}
 =2.04\times10^{-3}+2.96\times10^{-4}+4.27\times10^{-2}
 \approx0.0450~\centi\meter^{-1},
\]
and
\[
 \Sigma_{a,2}=1.91\times10^{-2}+1.38\times10^{-2}
 =0.0329~\centi\meter^{-1}.
\]
Therefore,
\[
 \phi_1\approx
 \frac{100}
 {0.0450-7.38\times10^{-4}
 -(3.36\times10^{-2})(4.27\times10^{-2})/0.0329}
 \approx1.45\times10^5~\centi\meter^{-2}\second^{-1},
\]
and
\[
 \phi_2=\frac{(4.27\times10^{-2})(1.45\times10^5)}{0.0329}
 \approx1.88\times10^5~\centi\meter^{-2}\second^{-1}.
\]
The volumetric absorption rate is
\[
\begin{aligned}
 R_a
 &=\phi_1\Sigma_{a,1}+\phi_2\Sigma_{a,2}\\
 &=1.45\times10^5(2.04\times10^{-3}+2.96\times10^{-4})
  +1.88\times10^5(0.0329)\\
 &\approx \boxed{6.5\times10^3~\second^{-1}\centi\meter^{-3}}.
\end{aligned}
\]
The volumetric fission rate is
\[
\begin{aligned}
 R_f
 &=\phi_1\Sigma_{f,1}+\phi_2\Sigma_{f,2}\\
 &=1.45\times10^5(2.96\times10^{-4})
  +1.88\times10^5(1.38\times10^{-2})\\
 &\approx \boxed{2.6\times10^3~\second^{-1}\centi\meter^{-3}}.
\end{aligned}
\]
\end{adaptedsolution}

\newpage

\noindent{\bf Problem 3}. Compute $k_{\infty}$ for the system described in
Problem 2.

\begin{adaptedsolution}
For the corresponding homogeneous eigenvalue problem,
\[
 k_\infty
 =\frac{\bar{\nu}\Sigma_{f,1}
 +\bar{\nu}\Sigma_{f,2}\Sigma_{s,2\gets1}/\Sigma_{a,2}}
 {\Sigma_{r,1}}.
\]
Substitution gives
\[
 k_\infty
 =\frac{7.38\times10^{-4}
 +(3.36\times10^{-2})(4.27\times10^{-2})/0.0329}
 {0.0450}
 \approx \boxed{0.985}.
\]
\end{adaptedsolution}

\clearpage
